\section{Conclusion}
In this paper, we present a practical kernel implementation of MoDiff to bridge the gap between theoretical quantization efficiency and real hardware acceleration for diffusion models. By introducing a cache-update fusion paradigm, our method eliminates additional I/O overhead during inference and enables efficient low-bit activation quantization in practice. Extensive experiments and ablation studies show that our implementation achieves up to 1.8x runtime speedup over FP32 models and reduces memory I/O usage by up to 42.2\%, demonstrating its effectiveness for accelerating diffusion model inference on hardware. These results highlight the promise of hardware-aware post-training quantization for improving the deployability of diffusion models in on-device and real-time scenarios. In future work, we will extend our evaluation to more datasets and model architectures, and compare against a broader set of quantization baselines to further validate the generality and effectiveness of our approach.

% Our empirical evaluation demonstrates that the proposed MoDiff implementation consistently delivers substantial and practically meaningful end-to-end performance improvements under realistic deployment conditions. The most significant gains arise in convolution-dominated workloads and large-batch inference settings, where low-bit kernels and fused execution effectively reduce both computational cost and memory traffic. While modest limitations remain due to quantization overhead and the presence of operators that cannot yet be quantized, the overall results clearly validate the practicality, robustness, and effectiveness of our approach. Taken together, these findings highlight the strong potential of MoDiff as a viable solution for efficient generative inference and reveal clear opportunities for further acceleration through deeper operator fusion, improved execution scheduling, and specialized kernel designs targeting emerging hardware architectures.